\section*{Conventions and manuscript concordance}
Every group is discrete unless we explicitly say profinite, pro-$\ell$, or
proalgebraic. All homomorphisms between profinite groups are continuous.
For a field $k$ of characteristic zero, a unipotent completion means the
universal prounipotent affine $k$-group, not a quotient of a profinite group
in the category of profinite groups. A genuine Galois action requires a
specified section or tangential base point; an outer action does not.
We write $\gamma_1 F=F$ and $\gamma_{r+1}F=[F,\gamma_rF]$.
A derivation of degree $m$ maps a generator of Lie degree $1$ to degree $m+1$.
This grading must not be confused with twice the motivic Tate weight.

\begin{longtable}{p{.19\textwidth}p{.38\textwidth}p{.34\textwidth}}
\toprule Our notation & Manuscript notation and meaning & Passage\\\midrule
$\Pi_0\subset\Pi$ & Discrete group and its specified profinite completion & \S52, f.750, scan205\\
$\mathcal G_0,\mathcal G$ & Discrete geometric extension and its closure; the closure need not be its full profinite completion & \S52, ff.750--751, scans205--207\\
$\mathcal M,\Gamma$ & Ambient profinite group and $\mathcal M/\mathcal G$ & \S52, ff.750--751, scans205--207\\
$\chi$ & Common cyclotomic character with open image & \S52, f.751, scan207\\
$\mathbb A_f$ & $\mathcal A_\Q=\widehat\Z\otimes\Q$ & \S52, f.748, scan201\\
$G,\rho$ & $\underline{\mathcal M}_\Q$, a finite-type algebraic target of an admissible adelic representation & \S52, ff.751--752, scans208--210\\
$\widetilde{\SL}_2(\Z)$ & Inverse image of $\SL_2(\Z)$ in the universal covering group of $\SL_2(\R)$ & \S53, ff.762--764, scans229--233\\
$z=(\sigma_1\sigma_2)^3$ & Central braid; the kernel over $\SL_2(\Z)$ is generated by $z^2$ & \S53, ff.764--766, scans233--237\\
$\operatorname{Cov}(S,\Phi)$ & Category of coverings of a germ; not a tower of Lie completions & \S53.5, ff.768--769, scans242--243\\
\bottomrule
\end{longtable}
Transcription marks, including doubtful and deleted readings, are evidence
to inspect, not axioms. The readings ledger records mathematical corrections
without claiming a final diplomatic reading of the handwriting.
